\section{Determinant of a Square Matrix}

Let n be a positive integer. Let us apply the above to the right vector space \(D_{d}^{n}\) over the field D; the elements of \(D_{d}^{n}\) are viewed as matrices with n rows and one column. Let \((\varepsilon_{1},\ldots,\varepsilon_{n})\) be the canonical basis of \(D_{d}^{n}\) . By Corollary 2 of VIII, p.~451, there exists a unique element \(\omega_{0}\) of \(\Omega(\mathrm{D}_{d}^{n})\) such that \(\omega_{0}(\varepsilon_{1},\ldots,\varepsilon_{n})=1\) . If A is an element of \(\mathbf{GL}_{n}(\mathrm{D})\) , then its columns \(a_{1},\ldots,a_{n}\) form a basis of \(D_{d}^{n}\) ; the element \(\omega_{0}(a_{1},\ldots,a_{n})\) of \(D_{ab}^{*}\) is called the determinant of A and denoted by \(\det(A)\) . Since we have \(a_{i}=A\varepsilon_{i}\) for \(1\leqslant i\leqslant n\) , the determinant of A is simply the determinant of the automorphism \(x\mapsto Ax\) of \(D_{d}^{n}\) . In particular, if the field D is commutative, then the determinant of A coincides with that defined in III, §8, No.~3, p.~524.

Let \(\mathrm{V}\) be a right vector space of finite dimension \(n\) over the field \(\mathrm{D}\) , and let \((e_1, \ldots, e_n)\) be a basis of \(\mathrm{V}\) . If \(u\) is an automorphism of \(\mathrm{V}\) and \(A\) is the matrix of \(u\) with respect to the basis \((e_1, \ldots, e_n)\) , then we have \(\det(u) = \det(A)\) .

We denote by \((E_{ij})\) the family of matrix units of \(\mathbf{M}_{n}(\mathrm{D})\) (II, §10, No.~3, p.~341). For every element \(\lambda\) of D and every pair \((i,j)\) of distinct integers in {[}1, n{]}, we set (II, §10, No.~13, p.~361)

\[
B _ {i j} (\lambda) = I _ {n} + \lambda E _ {i j};
\]

it is an element of \(\mathbf{GL}_{n}(\mathrm{D})\) . If \(\lambda_{1},\ldots,\lambda_{n}\) are elements of \(D^{*}\) , then the diagonal matrix \(\mathrm{diag}(\lambda_{1},\ldots,\lambda_{n})\) also belongs to \(\mathbf{GL}_{n}(\mathrm{D})\) .

PROPOSITION 3. --- The mapping det is the unique homomorphism from \(\mathbf{GL}_{n}(\mathrm{D})\) to \(D_{ab}^{*}\) that satisfies the relations

\[
\det (B _ {i j} (1)) = 1\tag{16}
\]

for \(i\neq j\) and

\[
\det (\mathrm{diag} (\lambda_ {1}, \ldots , \lambda_ {n})) = \pi (\lambda_ {1} \dots \lambda_ {n}),\tag{17}
\]

for \(\lambda_1, \ldots, \lambda_n \in \mathrm{D}^*\)

Let \(A\) be an element of \(\mathbf{GL}_n(\mathrm{D})\) and \(a_1, \ldots, a_n\) be its columns. We have

\[
\det (A) = \omega_ {0} (a _ {1}, \dots , a _ {n}).
\]

Now, the columns of the matrix \(A \operatorname{diag}(\lambda_1, \ldots, \lambda_n)\) are \(a_1\lambda_1, \ldots, a_n\lambda_n\) , and those of the matrix \(AB_{ij}(1)\) are therefore \(a_1, \ldots, a_j + a_i, a_{j+1}, \ldots, a_n\) . Since \(\omega_0\) is the unique element of \(\Omega(\mathrm{D}_d^n)\) such that \(\omega_0(\varepsilon_1, \ldots, \varepsilon_n) = 1\) , we see that the determinant is the unique mapping \(\varphi: \mathbf{GL}_n(\mathrm{D}) \to \mathrm{D}_{\mathrm{ab}}^*\) that satisfies the relations

\begin{enumerate}
\def\labelenumi{(\arabic{enumi})}
\setcounter{enumi}{17}
\tightlist
\item
\end{enumerate}

\[
\varphi (A \operatorname{diag} (\lambda_ {1}, \ldots , \lambda_ {n})) = \varphi (A) \pi (\lambda_ {1} \dots \lambda_ {n}),\tag{19}
\]

\[
\varphi (A B _ {i j} (1)) = \varphi (A) \quad \mathrm{for} i \neq j,\tag{20}
\]

\[
\varphi (I _ {n}) = 1.
\]

This proves, to begin with, that the determinant mapping satisfies relations (16) and (17). We already know that this mapping is a homomorphism from \(\mathbf{GL}_{n}(\mathrm{D})\) to \(D_{ab}^{*}\) . Conversely, if \(\varphi\) is a homomorphism from \(\mathbf{GL}_{n}(\mathrm{D})\) to \(D_{ab}^{*}\) such that \(\varphi(B_{ij}(1)) = 1\) for \(i \neq j\) and \(\varphi(\text{diag}(\lambda_1, \ldots, \lambda_n)) = \pi(\lambda_1 \cdots \lambda_n)\) , then \(\varphi\) satisfies relations (18) through (20) and is therefore equal to det.

Examples. --- 1) The columns of the matrix \(B_{ij}(\lambda)\) are \(\varepsilon_1, \ldots, \varepsilon_j + \varepsilon_i \lambda, \varepsilon_{i+1}, \ldots, \varepsilon_n\) . By formula (5) of VIII, p.~448, we have

\[
\det (B _ {i j} (\lambda)) = 1.\tag{21}
\]

\begin{enumerate}
\def\labelenumi{\arabic{enumi})}
\setcounter{enumi}{1}
\tightlist
\item
  Let \(\sigma\) be a permutation of the interval \([1,n]\) in \(\mathbf{N}\) , with signature \(\varepsilon(\sigma)\) . Let \(M(\sigma)\) be the matrix of the permutation \(\sigma\) (II, §10, No.~7, p.~351). The columns of the matrix \(M(\sigma)\) are \(\varepsilon_{\sigma(1)},\ldots,\varepsilon_{\sigma(n)}\) . Applying formula (4) of VIII, p.~448, we therefore have
\end{enumerate}

\[
\det (M (\sigma)) = \pi (\varepsilon (\sigma)).\tag{22}
\]

\begin{enumerate}
\def\labelenumi{\arabic{enumi})}
\setcounter{enumi}{2}
\item
  Suppose \(n \geqslant 1\) . For every invertible diagonal matrix of the form \(\Delta = \text{diag}(d_1, \ldots, d_n)\) , we have \(\Delta B_{ij}(\lambda)\Delta^{-1} = B_{ij}(d_i\lambda d_j^{-1})\) . Let A be an element of \(\mathbf{GL}_n(\mathrm{D})\) . By Corollary 1 of II, §10, No.~13, p.~362 and the previous formula, there exist matrices P and \(\Delta\) in \(\mathbf{GL}_n(\mathrm{D})\) such that \(A = P\Delta\) , that P is the product of matrices of the form \(B_{ij}(\lambda)\) , and that \(\Delta\) is a diagonal matrix of the form \(\text{diag}(1, \ldots, 1, d)\) . We have \(\det(P) = 1\) by Example 1, and therefore \(\det(A) = \det(\Delta) = \pi(d)\) by Proposition 3.
\item
  Let \(D'\) be a field, and let u be a homomorphism from D to \(D'\) . When passing to the quotients, u defines a group homomorphism \(u_{ab}\) from \(D_{ab}^{*}\) to \(D_{ab}^{\prime *}\) . Let \(u_{n}\) be the homomorphism from \(\mathbf{GL}_{n}(\mathrm{D})\) to \(\mathbf{GL}_{n}(\mathrm{D}')\) that transforms a matrix \(A = (a_{ij})\) into the matrix \((u(a_{ij}))\) . The formula
\end{enumerate}

\[
\det (u _ {n} (A)) = u _ {\mathrm{ab}} (\det (A))\tag{23}
\]

for \(A \in \mathbf{GL}_n(\mathrm{D})\) follows immediately from Example 3.

\begin{enumerate}
\def\labelenumi{\arabic{enumi})}
\setcounter{enumi}{4}
\item
  We have \(\mathbf{GL}_1(\mathrm{D}) = \mathrm{D}^*\) , and Proposition 3 shows that we have \(\det(a) = \pi(a)\) for \(a\) in \(\mathbf{GL}_1(\mathrm{D})\) .
\item
  Suppose \(n = 2\) . Let \(A = \begin{pmatrix} a & b \\ c & d \end{pmatrix}\) be an element of \(\mathbf{GL}_2(\mathrm{D})\) . The elements \(a\) and \(c\) of \(\mathrm{D}\) are not both zero. We will give the determinant of \(A\) explicitly.
\end{enumerate}

If \(a\) is not zero, then we have

\[
A = \left( \begin{array}{c c} 1 & 0 \\ c a ^ {- 1} & 1 \end{array} \right) \left( \begin{array}{c c} a & 0 \\ 0 & d - c a ^ {- 1} b \end{array} \right) \left( \begin{array}{c c} 1 & a ^ {- 1} b \\ 0 & 1 \end{array} \right).\tag{24}
\]

Consequently, \(ad - aca^{-1}b \neq 0\) and

\[
\det (A) = \pi (a d - a c a ^ {- 1} b).\tag{25}
\]

\begin{enumerate}
\def\labelenumi{\alph{enumi})}
\setcounter{enumi}{1}
\tightlist
\item
  If \(a\) is zero, then we have \(c \neq 0\) and
\end{enumerate}

\[
A = \left( \begin{array}{c c} 0 & b \\ c & d \end{array} \right) = \left( \begin{array}{c c} 0 & 1 \\ 1 & 0 \end{array} \right) \left( \begin{array}{c c} c & d \\ 0 & b \end{array} \right).\tag{26}
\]

Consequently, by a) and Example 2, we have \(cb \neq 0\) and

\[
\det (A) = \pi (- c b).\tag{27}
\]

\begin{enumerate}
\def\labelenumi{\arabic{enumi})}
\setcounter{enumi}{6}
\tightlist
\item
  Let \(D^{o}\) be the opposite field of D. The mapping \(A \mapsto {}^{t}A\) from \(\mathbf{M}_{n}(\mathrm{D})\) to \(\mathbf{M}_{n}(\mathrm{D}^{\circ})\) satisfies \({}^{t}(AB) = {}^{t}B^{t}A\) . Therefore, for every matrix A in \(\mathbf{GL}_{n}(\mathrm{D})\) , the transposed matrix \({}^{t}A\) is invertible in \(\mathbf{M}_{n}(\mathrm{D}^{\circ})\) , but it is not necessarily invertible in \(\mathbf{M}_{n}(\mathrm{D})\) . It follows from Example 3 that if A belongs to \(\mathbf{GL}_{n}(\mathrm{D})\) , then the elements A of \(\mathbf{GL}_{n}(\mathrm{D})\) and \({}^{t}A\) of \(\mathbf{GL}_{n}(\mathrm{D}^{\circ})\) have the same determinant. On the other hand, even if the matrix \({}^{t}A\) belongs to \(\mathbf{GL}_{n}(\mathrm{D})\) , its determinant in \(\mathrm{GL}_{n}(\mathrm{D})\) is not necessarily equal to that of A (cf.~Example 6).
\end{enumerate}

Remarks. --- 1) Let V be a right vector space over the field D, of finite dimension n.~We view the dual space \(V^{*} = \operatorname{Hom}_{\mathrm{D}}(\mathrm{V}, \mathrm{D})\) as a right vector space over the opposite field \(D^{o}\) of D. Let u be an automorphism of V, and let \(^{t}u\) be the automorphism of \(V^{*}\) that is the transpose of u. If u is represented by a matrix A in \(\mathbf{M}_{n}(\mathrm{D})\) with respect to a basis \((e_{1}, \ldots, e_{n})\) of V, then the automorphism \(^{t}u\) is represented by the matrix \(^{t}A\) in \(\mathbf{M}_{n}(\mathrm{D}^{o})\) with respect to the basis \((e_{1}^{*}, \ldots, e_{n}^{*})\) of \(V^{*}\) dual to \((e_{1}, \ldots, e_{n})\) (II, §10, No.~4, p.~344, Proposition 3). By Example 7, the determinant of \(^{t}u\) is equal to that of u.

\begin{enumerate}
\def\labelenumi{\arabic{enumi})}
\setcounter{enumi}{1}
\tightlist
\item
  The results of Subsections 1 through 4 generalize to the case when D is a local ring (VIII, p.~459, Exercise 2).
\end{enumerate}

